\documentclass[11pt]{article}
\usepackage[T1]{fontenc}
\usepackage{lmodern}
\usepackage{amsmath,amssymb,amsthm,mathtools}
\usepackage{graphicx}
\usepackage[margin=1in]{geometry}
\usepackage{microtype}
\usepackage{enumitem}
\usepackage[hidelinks]{hyperref}

\newtheorem{theorem}{Theorem}
\newtheorem{lemma}[theorem]{Lemma}
\newtheorem{corollary}[theorem]{Corollary}
\newcommand{\A}{\mathcal A}
\newcommand{\D}{\mathcal D}
\newcommand{\Prim}{\operatorname{Prim}}
\newcommand{\hull}{\operatorname{hull}}
\newcommand{\Ad}{\operatorname{Ad}}

\title{Finite Primitive Sets Are Sets of Synthesis for\\
$p$-Adic Unipotent Weighted Orlicz Algebras}
\author{Automated Mathematics Run \texttt{fa\_banach\_001}, lane 2}
\date{21 July 2026}

\begin{document}
\maketitle

\begin{abstract}
Question 8.12 of Carter's paper asks for geometric conditions on subsets of
the primitive and coadjoint-orbit spectra which imply spectral synthesis for
Hermitian weighted Orlicz algebras on unipotent $p$-adic groups.  We prove
that every finite subset is a set of synthesis.  More generally, the result
holds for every finite closed family of admissible primitive points of a
locally elliptic group.  Carter's CCR theorem makes all primitive points
admissible in the $p$-adic unipotent case.  The proof combines compact-open
smoothing, finite-dimensional representation corners, finite Hecke algebras,
and an exact central-support cutoff in Carter's generating ideal $j(C)$.
We also retain the earlier compact-open synthesis theorem in the abelian case.
\end{abstract}

\paragraph{Status.}
Candidate substantial partial result, likely valid, subject to human review.
It completely establishes the geometric criterion ``$C$ is finite'' for all
groups in Question 8.12, including nonabelian ones.  The source question is
open-ended, and infinite subsets of nonabelian orbit spaces remain untreated.

\section{Source question}

The source is Max Carter, \emph{Weighted Orlicz $*$-algebras on locally
elliptic groups}, arXiv:2503.20735v2, Question 8.12 on page 29
\cite{Carter}.  Given a $p$-adic unipotent matrix group $N$, Carter identifies
\[
 \Prim_*(L^\Phi(N,\omega))\cong
 \mathfrak n^*/\Ad^*(N)
\]
and asks which geometric properties of a subset $C$ imply
$\ker(C)=j(C)$.

\begin{figure}[ht]
  \centering
  \includegraphics[width=\textwidth]{figures/open_problem_crop.png}
  \caption{Question 8.12, rendered from page 29 of
  arXiv:2503.20735v2.}
\end{figure}

\section{Definitions}

Let $G$ be locally elliptic, let $(\Phi,\Psi)$ be a complementary pair of
Young functions with $\Phi\in\Delta_2$, and let $\omega$ be an
$L^\Phi$-weight with $1/\omega\in L^\Psi(G)$.  Put
\[
  \A=L^\Phi(G,\omega)
\]
and assume that $\A$ is Hermitian.  Carter identifies the irreducible
$*$-representations of $\A$ with the irreducible unitary representations of
$G$ and defines, for closed $C\subseteq\Prim_*(\A)$,
\[
 \ker(C)=\bigcap_{P\in C}P.
\]
His smooth functional calculus produces a closed two-sided ideal $j(C)$ with
$j(C)\subseteq\ker(C)$ and $\hull_*(j(C))=C$; the set $C$ is a
\emph{set of synthesis} if $j(C)=\ker(C)$ \cite[Theorem 7.1]{Carter}.

Write $\D(G)$ for the compactly supported locally constant functions on $G$.
If $H\leq G$ is compact-open, set
\[
 e_H=\mu(H)^{-1}1_H.
\]
Then $e_H=e_H^*=e_H*e_H$, and $\pi(e_H)$ is the orthogonal projection onto
the $H$-fixed vectors of a unitary representation $\pi$.  Call $\pi$
\emph{admissible} if these fixed-vector spaces are finite dimensional for
all compact-open $H$.

\section{Proof intuition}

There are two scales in the proof.  At the analytic scale, convolution on
both sides by $e_H$ approximates the identity.  For a finite family of
admissible representations, the smoothed algebra $e_H\A e_H$ is evaluated
inside a finite direct sum of matrix algebras.  A bounded linear lift in that
finite-dimensional image corrects dense test functions so that all selected
representations annihilate them exactly.

At the algebraic scale, every test function lies in a finite-dimensional
Hecke algebra.  If all selected representations annihilate a test function
$d$, they also annihilate the central support projection $p_d$ of $d$ in
that Hecke algebra.  Carter's smooth functional calculus turns a scalar
multiple of $p_d$ back into exactly $p_d$, placing it in $m(C)$.  Since
$d=p_d*d$, the test function lies in $j(C)$.  Passing first through the
finite corner and then through the compact-open approximate identity proves
synthesis.  In the $p$-adic unipotent case, CCR makes every fixed-vector
projection compact and therefore finite rank.

\section{Finite Hecke support}

\begin{lemma}[Dense locally finite test algebra]\label{lem:test}
$\D(G)$ is dense in $\A$, and every $d\in\D(G)$ belongs to a
finite-dimensional convolution $*$-subalgebra of $\D(G)$.
\end{lemma}

\begin{proof}
The $\Delta_2$ hypothesis gives density of compactly supported functions in
the weighted Orlicz space.  Since a locally elliptic group has a compact-open
basis, uniform approximation on compact sets by locally constant step
functions, together with local boundedness of $\omega$, gives density of
$\D(G)$.

Fix $d\in\D(G)$.  Its compact support is contained in a compact-open subgroup
$K$.  Uniform local constancy on $K$ supplies an open subgroup of $K$ under
whose left and right translations $d$ is invariant.  Replacing it by its
normal core gives an open normal subgroup $H\triangleleft K$ with the same
property.  The $H$-bi-invariant functions supported in $K$ form a
finite-dimensional convolution $*$-algebra, isomorphic after Haar
normalization to the finite group algebra $\mathbb C[K/H]$, and contain $d$.
\end{proof}

\begin{lemma}[Algebraic kernel elements lie in $j(C)$]\label{lem:algebraic}
Let $C\subseteq\Prim_*(\A)$ be closed.  Then
\[
  \D(G)\cap\ker(C)\subseteq j(C).
\]
\end{lemma}

\begin{proof}
Take $d\in\D(G)\cap\ker(C)$ and place it, using
Lemma~\ref{lem:test}, in a finite-dimensional convolution $*$-algebra $B$.
Viewed in its faithful finite-dimensional regular representation, $B$ is a
finite-dimensional $C^*$-algebra.  Let $p_d$ be the central support projection
of $d$ in $B$, so the ideal generated by $d$ is $Bp_d$ and $p_d*d=d$.

For every irreducible representation $\pi$ whose kernel lies in $C$, the
kernel of the restricted $*$-representation $\pi|_B$ is an ideal containing
$d$.  It therefore contains $p_d$, so $\pi(p_d)=0$.  Choose
\[
 0<a<2\pi,\qquad a\lVert p_d\rVert_1\leq1.
\]
Carter's cutoff Proposition 5.5 provides $\varphi\in A_\gamma$ which is zero
near $0$ and satisfies $\varphi(a)=1$.  Thus $\varphi\{ap_d\}\in m(C)$.
Since $p_d*p_d=p_d$ and $\varphi(0)=0$, the defining power series for the
functional calculus gives
\[
 \varphi\{ap_d\}=\varphi(a)p_d=p_d.
\]
Hence $p_d\in m(C)\subseteq j(C)$, and $d=p_d*d\in j(C)$.
\end{proof}

\section{Finite admissible sets}

\begin{theorem}[Finite admissible synthesis]\label{thm:admissible}
Let $C=\{\ker\pi_1,\ldots,\ker\pi_r\}$ be a finite closed subset of
$\Prim_*(\A)$.  If every $\pi_i$ is admissible, then $C$ is a set of
synthesis.
\end{theorem}

\begin{proof}
Put $I=\ker(C)$.  Fix a compact-open subgroup $H\leq G$ and define
\[
 \A_H=e_H*\A*e_H,
 \qquad
 \D_H=e_H*\D(G)*e_H.
\]
The density in Lemma~\ref{lem:test} and boundedness of convolution by $e_H$
show that $\D_H$ is dense in $\A_H$.  Consider
\[
 T_H:\A_H\longrightarrow
 \bigoplus_{i=1}^r
 \mathcal B(\mathcal H_{\pi_i}^H),
 \qquad T_H(f)=(\pi_1(f),\ldots,\pi_r(f)).
\]
The codomain is finite dimensional by admissibility.  Therefore
$T_H(\D_H)$ is a closed linear subspace.  Density implies
\[
 T_H(\D_H)=T_H(\A_H).
\]
Choose a linear right inverse
\[
 R_H:T_H(\A_H)\longrightarrow\D_H
\]
for $T_H$.  Its domain is finite dimensional, so $R_H$ is bounded for the
$\A$-norm on its range.

Let $f\in I$ and put $f_H=e_H*f*e_H$.  Then $f_H\in I\cap\A_H$.  Choose
$d_n\in\D_H$ with $d_n\to f_H$ in $\A$ and set
\[
 h_n=d_n-R_H(T_H(d_n)).
\]
Since $T_H(f_H)=0$, continuity gives $T_H(d_n)\to0$, and boundedness of $R_H$
gives $h_n\to f_H$.  Moreover $T_H(h_n)=0$, so
$h_n\in\D(G)\cap I$.  Lemma~\ref{lem:algebraic} yields $h_n\in j(C)$, and
closedness gives $f_H\in j(C)$.

As $H$ shrinks to the identity, Carter's compact-neighbourhood approximate
identity \cite[Proposition 2.16(ii)]{Carter} gives
\[
 e_H*f*e_H\longrightarrow f
 \quad\text{in }\A.
\]
Thus $f\in j(C)$.  We have proved $I\subseteq j(C)$; the reverse inclusion is
part of Carter's construction, so $I=j(C)$.
\end{proof}

\begin{corollary}[All finite coadjoint-orbit sets]\label{cor:padic}
Let $N$ be a group of unipotent matrices over a $p$-adic field, and let
$\A=L^\Phi(N,\omega)$ satisfy the hypotheses above.  Then every finite subset
of
\[
 \mathfrak n^*/\Ad^*(N)\cong\Prim_*(\A)
\]
is a set of synthesis.
\end{corollary}

\begin{proof}
Carter's Theorem 8.10 states that $N$ is CCR and identifies its unitary dual
with the coadjoint-orbit space.  If $H\leq N$ is compact-open and $\pi$ is
irreducible, then $\pi(e_H)$ belongs to
$\pi(C^*(N))=\mathcal K(\mathcal H_\pi)$.  It is also an orthogonal
projection, hence has finite rank.  Thus every irreducible representation is
admissible.  The CCR property also makes primitive ideals maximal, so the
primitive spectrum is $T_1$ and every finite subset is closed.  Apply
Theorem~\ref{thm:admissible} and Carter's Corollary 8.11.
\end{proof}

\section{The retained abelian compact-open criterion}

\begin{theorem}\label{thm:compactopen}
If $G$ is abelian, every compact-open subset of the Pontryagin dual
$\widehat G\cong\Prim_*(\A)$ is a set of synthesis.
\end{theorem}

\begin{proof}
Fourier transformation maps $\D(G)$ onto $\D(\widehat G)$.  If a compact-open
$U\subseteq\widehat G$ is disjoint from a closed $C$, let
$e_U\in\D(G)$ satisfy $\widehat e_U=1_U$.  Scale $e_U$ to $h=ae_U$ with
$0<a<2\pi$ and $\lVert h\rVert_1\leq1$.  Every representation in $C$
annihilates $h$.  A Carter cutoff which is zero near $0$ and one at $a$ gives
$\varphi\{h\}=e_U\in m(C)$.  Hence every test function whose Fourier support
is contained in $U$ belongs to $j(C)$.

Now let $C$ itself be compact-open and choose $q_C\in\D(G)$ with
$\widehat q_C=1_C$.  For $f\in\ker(C)$, Fourier injectivity gives
$f*q_C=0$.  If $f_n\in\D(G)$ tends to $f$, put
\[
 g_n=f_n-f_n*q_C.
\]
Then
\[
 \lVert g_n-f\rVert_\A
 \leq(1+\lVert q_C\rVert_\A)\lVert f_n-f\rVert_\A\longrightarrow0,
\]
and the compact-open support of $\widehat g_n$ is disjoint from $C$.
The cutoff argument gives $g_n\in j(C)$, hence $f\in j(C)$.  Therefore
$\ker(C)=j(C)$.
\end{proof}

\section{Verification, limitations, and novelty}

The proof uses no numerical computation.  Its critical points were checked
independently: compact-open indicators are the source's approximate identity;
CCR turns their represented projections into finite-rank projections; the
right inverse is used only after fixing $H$, so no uniform bound as
$H\downarrow\{e\}$ is required; and the central-support projection belongs
to Carter's exact generating set $m(C)$.

The result does not classify infinite subsets of a nonabelian orbit space.
Local indexes and bounded arXiv searches on 21 July 2026 used the exact source
id and combinations of ``spectral synthesis'', ``$p$-adic unipotent'',
``finite sets'', ``singleton'', ``primitive spectrum'', and ``weighted
Orlicz''.  No exact prior theorem was found.  Carter explicitly notes that
even the unweighted $L^1(N)$ case would be novel \cite[Section 8.3]{Carter};
novelty confidence is nevertheless only moderate pending expert review.

\paragraph{Human-review recommendation.}
Check Lemma~\ref{lem:algebraic}, especially the central support inside the
finite Hecke algebra, and check the equality
$T_H(\D_H)=T_H(\A_H)$ before the finite-dimensional lift is chosen.  These
are the only new structural steps.

\begin{thebibliography}{9}
\bibitem{Carter}
M.~Carter,
\emph{Weighted Orlicz $*$-algebras on locally elliptic groups},
Studia Math. \textbf{287} (2026), no.~2, 145--180;
arXiv:2503.20735v2.
\end{thebibliography}

\end{document}
